\section{Adaptive-Agent Mechanism}

A central modeling question is whether positive outcome--accessibility alignment must be inserted exogenously. We instead consider an adapted agent with environment state $X_t$, internal belief state $B_t$, scalar evaluation signal $Y_t=h(B_t)$, and later outcome $U_T$.

Adaptation can create statistical dependence between the internal model and future outcomes, but mutual information alone is not enough for the mean-uplift result. The relevant directional object is the conditional mean
\[
m(y)=\mathbb E[U_T\mid Y_t=y].
\]
Suppose accessibility is score-measurable,
\[
S_t=s(Y_t),
\]
with $S_t\ge0$, $0<\mathbb E[S_t]<\infty$, and the required product integrability. Then Supplementary Theorem S2 gives the exact projection identity
\[
\operatorname{Cov}(U_T,S_t)
=
\operatorname{Cov}(m(Y_t),s(Y_t)).
\]
If $m(y)$ and $s(y)$ are both nondecreasing on the support of $Y_t$, then
\[
\operatorname{Cov}(U_T,S_t)\ge0.
\]
The inequality is strict when the two orderings vary together on a positive-probability set. By the main covariance identity, this implies nonnegative first-person mean uplift under the corresponding weighted-measure model.

A particularly transparent corollary is
\[
S_t=r(m(Y_t))
\]
for a nondecreasing function $r$. If the predicted conditional mean is nonconstant and $r$ is strictly increasing on its essential range, the covariance is strictly positive.

There is also an exact self-calibration case. If the agent's score is its posterior-mean forecast from internal information $B_t$,
\[
Y_t=\mathbb E[U_T\mid B_t],
\]
then the tower property gives
\[
\mathbb E[U_T\mid Y_t]
=
\mathbb E[\mathbb E[U_T\mid B_t]\mid Y_t]
=
Y_t.
\]
Thus the directional-calibration premise of S2 is automatic for a true posterior-mean score. Any nondecreasing accessibility function $s(Y_t)$ then gives nonnegative covariance, with strict positivity when the score is nonconstant and the accessibility map is strictly increasing on its essential range.

For a learned score that is only approximately calibrated, let
\[
e(Y_t)=\mathbb E[U_T\mid Y_t]-Y_t.
\]
Under square integrability,
\[
\operatorname{Cov}(U_T,S_t)
=
\operatorname{Cov}(Y_t,S_t)
+
\operatorname{Cov}(e(Y_t),S_t),
\]
and Cauchy--Schwarz yields
\[
\operatorname{Cov}(U_T,S_t)
\ge
\operatorname{Cov}(Y_t,S_t)
-
\sqrt{\operatorname{Var}(e(Y_t))\operatorname{Var}(S_t)}.
\]
Hence a sufficient robustness condition is
\[
\operatorname{Cov}(Y_t,S_t)
>
\sqrt{\operatorname{Var}(e(Y_t))\operatorname{Var}(S_t)}.
\]

The calibration-error term can itself be bounded by ordinary prediction mean-squared error. Since
\[
e(Y_t)=\mathbb E[U_T-Y_t\mid Y_t],
\]
conditional Jensen implies
\[
\operatorname{Var}(e(Y_t))
\le
\mathbb E[e(Y_t)^2]
\le
\mathbb E[(U_T-Y_t)^2].
\]
Therefore
\[
\operatorname{Cov}(U_T,S_t)
\ge
\operatorname{Cov}(Y_t,S_t)
-
\sqrt{\mathbb E[(U_T-Y_t)^2]\operatorname{Var}(S_t)}.
\]
This S2.4 population certificate is conservative because prediction MSE also includes irreducible conditional outcome variance.

Corollary S2.5 converts the population inequality into a bounded independent-held-out certificate. Corollary S2.6 shows that arbitrary upstream training is compatible with this guarantee when the final certification sample is independent. Corollary S2.7 permits same-holdout selection among a finite predeclared candidate family after multiplicity correction.

Corollary S2.8 then separates the QBS covariance-composition step from any particular concentration inequality. It requires only a simultaneous confidence envelope for
\[
\mathbb E[Y_t],
\quad
\mathbb E[S_t],
\quad
\mathbb E[Y_tS_t],
\quad
\mathbb E[S_t^2],
\quad
\mathbb E[(U_T-Y_t)^2].
\]
Any statistical method that validly supplies those five simultaneous moment bounds can be inserted without changing the QBS covariance proof.

Corollary S2.9 supplies one unbounded light-tail instantiation. It uses sub-Gaussian sample-mean control for $Y_t$ and $S_t$ and explicit Bernstein/sub-exponential sample-mean controls for
\[
Y_tS_t,
\qquad
S_t^2,
\qquad
(U_T-Y_t)^2.
\]
With
\[
t_\delta=\log\frac{10}{\delta},
\]
the five concentration events form a simultaneous S2.8 envelope with probability at least $1-\delta$. The resulting lower margin $D_{\rm LT}$ therefore satisfies
\[
P\!\left(
\operatorname{Cov}(U_T,S_t)\ge D_{\rm LT}
\right)
\ge1-\delta.
\]
This permits unbounded light-tail validation without silently assuming universal product or square constants from marginal sub-Gaussianity.

For square-integrable $U_T$, the quantity
\[
M(U_T;Y_t)
=
\operatorname{Var}(\mathbb E[U_T\mid Y_t])
\]
measures conditional-mean predictive strength. Positive mutual information $I(U_T;Y_t)>0$ is weaker: dependence can occur entirely through conditional variance while $\mathbb E[U_T\mid Y_t]$ remains constant, in which case every score-measurable accessibility map has zero covariance with $U_T$.

The intended causal and validation chain is therefore
\[
\text{adaptation / inference}
\rightarrow
\text{posterior-mean or approximately calibrated score}
\rightarrow
\text{population covariance certificate}
\rightarrow
\text{finite-sample moment envelopes}
\rightarrow
\text{selection-safe statistical certification}.
\]

Experiment E2 tests the minimal version of the predictive mechanism in a nonlinear environment. A misspecified model that cannot represent the relevant interaction fails to produce predictive alignment, while a small interaction-capable model succeeds. E3 supplies a paired endogenous-policy example. These simulations support premise generation in toy environments; S2--S2.9 provide the mathematical and statistical implications once the stated calibration, accessibility, concentration, held-out-evaluation, and selection conditions are specified. None of these steps establishes a physical Everettian accessibility map.
